\chapter{Deployment and Conclusion}

The final phase of the software development life cycle is the deployment and delivery of the system to the production environment. This chapter discusses the deployment strategy used for the JUST Undergraduate Final Admission System (JUST-UFAS) and the future prospects of the project.

\section{Deployment Process}
Deploying a modern web application built with Next.js and Node.js involves a structured workflow to ensure stability and performance. The deployment process for JUST-UFAS followed these essential steps:

\begin{itemize}
    \item \textbf{Production Build:} The application code was optimized and compiled into a production-ready build to ensure fast loading times and efficient resource management.
    \item \textbf{Environment Configuration:} Essential environment variables, such as PostgreSQL database connection strings and Prisma client settings, were configured on the production server.
    \item \textbf{Server Setup:} The application is hosted on the university's official infrastructure. Due to existing resource allocations, the system is currently deployed on a specific port.
    \item \textbf{Maintenance and Monitoring:} The system is actively maintained by the \textbf{Technical Committee of Undergraduate Admission} at Jashore University of Science and Technology.
\end{itemize}

\section{Deployment Status and URL}
The system is currently live and accessible for the final admission process. 

\begin{itemize}
    \item \textbf{Current Hosting:} The project is running on the university's existing domain using port 3000 as a cost-effective hosting solution.
\end{itemize}

\section{Conclusion}
The "JUST Undergraduate Final Admission System" offers a significant improvement over manual admission processes. By automating the workflow, the system has achieved:
\begin{itemize}
    \item \textbf{Efficiency:} Faster data processing and reduced manual labor.
    \item \textbf{Accuracy:} Minimized human errors in student data handling and merit list verification.
    \item \textbf{Cost-Effectiveness:} Reduced need for extra manpower, leading to operational cost savings for the university.
\end{itemize}

\section{Future Work}
While the current system fulfills all primary requirements, future updates could include:
\begin{itemize}
    \item \textbf{Automated Migration:} Implementing a full auto-migration system for students across different departments.
    \item \textbf{Payment Gateway Integration:} Integrating local mobile banking (bKash, Nagad) for seamless admission fee collection.
    \item \textbf{Mobile Application:} Developing a dedicated mobile app for students to track their admission status in real-time.
\end{itemize}

Through this project, Jashore University of Science and Technology has taken a step forward in digitizing its administrative infrastructure, ensuring a smoother experience for future students.